\documentclass[12pt,uplatex]{jsarticle}   % 日本語
\usepackage{tcu-thesis}
\usepackage{fancyhdr}
\usepackage{bm} % 太字の数式
\usepackage{amssymb} % 数式
\usepackage{amsmath} % 数式
\usepackage{float} % 図表の位置
\usepackage{listings} % ソースコード
\usepackage{color} % ソースコードの色付け

% ソースコードのスタイル設定
\lstset{
  basicstyle=\ttfamily\footnotesize,
  numbers=left,
  numberstyle=\tiny\color{gray},
  stepnumber=1,
  numbersep=5pt,
  backgroundcolor=\color{white},
  showspaces=false,
  showstringspaces=false,
  showtabs=false,
  frame=single,
  rulecolor=\color{black},
  tabsize=2,
  captionpos=b,
  breaklines=true,
  breakatwhitespace=false,
  title=\lstname,
  keywordstyle=\color{blue},
  commentstyle=\color{green},
  stringstyle=\color{red},
  escapeinside={\%*}{*)},
  morekeywords={*,...}
}

% Python用のスタイル設定
\lstdefinestyle{python}{
  language=Python,
  keywordstyle=\color{blue},
  commentstyle=\color[RGB]{0,100,0},     %DarkGreen
  stringstyle=\color[RGB]{128,0,0},      %Maroon
  frame=single,
  breaklines=true,
  basicstyle=\ttfamily\footnotesize
}

% Jupyter Notebook用のスタイル設定
\lstdefinestyle{jupyter}{
  language=Python,
  keywordstyle=\color{blue},
  commentstyle=\color[RGB]{0,100,0},     %DarkGreen
  stringstyle=\color[RGB]{128,0,0},      %Maroon
  frame=single,
  breaklines=true,
  basicstyle=\ttfamily\footnotesize,
  morekeywords={In, Out, print}
}


% アルゴリズムのレイアウト
\newenvironment{algorithm_step}
{\begin{list}{}
    {\setlength{\leftmargin}{6em}  % Step X)の部分の幅
     \setlength{\itemindent}{0em}
     \setlength{\labelsep}{0em}
     \setlength{\labelwidth}{6em}  % Step X)の部分の幅
     \setlength{\itemsep}{1em}     % ステップ間の縦方向の間隔
     \renewcommand{\makelabel}[1]{##1\hfill}}}
{\end{list}}

% 数式カウンターの定義
\newcounter{myeqcounter} % 独自カウンターを定義

\newcommand{\autoequation}[1]{
\refstepcounter{myeqcounter} % カウンターを進める
\begin{equation}
#1 \tag{\themyeqcounter} % カウンター値をタグに適用
\end{equation}
}


% オリジナルマクロ by Kyohei Fushida and Kenji Fujiwara
\newcommand{\li}{\item}
\newcommand{\ol}[1]{\begin{enumerate}#1\end{enumerate}}
\newcommand{\ul}[1]{\begin{itemize}#1\end{itemize}}
\newcommand{\dl}[1]{\begin{description}#1\end{description}}
\newcommand{\equ}[1]{\begin{equation}#1\end{equation}}
\newcommand{\eqenum}[1]{\begin{align}#1\end{align}}
\newcommand{\eqa}[1]{\begin{align}#1\end{align}}
\newcommand{\eqmul}[1]{\begin{equation}\begin{split}#1\end{split}\end{equation}}
\usepackage{xcolor}

%\newcommand{\figref}[1]{図\ref{#1}}
%\newcommand{\tabref}[1]{表\ref{#1}}
\renewcommand{\subsecref}[1]{\ref{#1}項}
\newcommand{\chapref}[1]{\ref{#1}章}

\definecolor{darkgreen}{rgb}{0, 0.5, 0} %下の色指定で使う
\definecolor{whitesmoke}{rgb}{0.99, 0.99, 0.99} %下の色指定で使う
\newcommand{\red}{\color{red}}%red
\newcommand{\blue}{\color{blue}}
\newcommand{\green}{\color{teal}}
\newcommand{\purple}{\color{violet}}%violet
\newcommand{\black}{\color{black}}
\newcommand{\ea}{\xspace\emph{et al.}\xspace}
\newcommand{\TODO}[1]{\nb{TODO:}{#1}}


\newcommand{\methods}{\#methods}
\newcommand{\colorrow}{\rowcolor[rgb]{0.851,0.851,0.851}}
\newcommand{\colorcel}{\cellcolor[rgb]{0.851,0.851,0.851}}
\def\et{\xspace et\ al.\xspace}
\def\ie{i.e.,\xspace}
\def\eg{e.g.,\xspace}
\newcommand{\nb}[2]{
    \fcolorbox{gray}{yellow}{\bfseries\sffamily\scriptsize#1}
    {\sf\small\textit{\textcolor{blue}{#2}}}
   }

\newtheorem{defi}{定義}
\newtheorem{hypo}{仮説}

% 必要なパッケージの読み込み
% とりあえず最低限使っておいて欲しいパッケージのみ書いています．
% 色々なパッケージがあります．藤原が過去に使ってたパッケージなどは博士論文のファイルを参照してください．
\usepackage[dvipdfmx]{graphicx}
\usepackage{cite} % 参考文献を [3,4]のようにまとめるのに便利．
\usepackage{url} % urlを記載するときに使う．

% 参考文献を引用された順で出力する
\bibliographystyle{junsrt}

% 学籍番号
\studentnumber{20XXXXXX}

% 日本語題目
\title{論文タイトル}

% 英語題目
\etitle{Thesis Title}

% 日本語氏名（姓と名の間に空白（半角）を入れて下さい）
\author{姓 名}

% 欧文氏名（first name, last name の順に記入し，先頭文字のみを大文字にする．）
\eauthor{First Last}

% 専攻
\department{\media}
\course{\jsys}

% 論文提出年月日
\syear{20XX}
\smonth{1}
\sday{1}

% % キーワード５〜６個 (in LaTeX)
% \keywords{キーワード1，キーワード2，キーワード3}

% % 5 or 6 Keywords (in LaTeX)
% \ekeywords{Keyword 1, Keyword 2, Keyword 3}

% 内容梗概 (in LaTeX)
%   Abstractは必須です．
%   サンプルほど長々と書く必要はありません．
%   注: 行の先頭が\\で始まらないようにすること．
% \abstract{%
% ここに内容梗概を記述する．
% 研究背景，研究目的，提案内容，実験結果，結論を簡潔にまとめる．
% 行の先頭が \\ で始まらないように注意する．
% }

%%%%%%%%%%%%%%%%%%%%%%%%% document starts here %%%%%%%%%%%%%%%%%%%%%%%%%%%%

\begin{document}
%
% 表紙 および アブストラクト
%

\begin{center}
    % \vspace*{3cm} % 上部の余白調整

    {\Large 20XX 年度 卒業論文}
    \vspace{3cm}

    {\LARGE 論文タイトル}
    \vspace{5cm}

    \begin{flushright}
        {\large 東京都市大学} \\
        {\large メディア情報学部情報システム学科} \\
        {\large 三川研究室} \\
        {\large 2XXXXXX 姓 名} \\
        {\large 指導: 三川 健太 准教授} \\
    \end{flushright}

\end{center}

\thispagestyle{empty}
% \titlepage
% \jabstractpage
%
% 目次
%
\clearpage
\pagenumbering{roman} % ページ番号をローマ数字にする
\tableofcontents
\clearpage
\newpage

% 以降本文
\pagenumbering{arabic} % ページ番号をアラビア数字に戻す

% この章構成はあくまで一例です．
% \section:章
% \subsection: 節
% \subsubsection: 項
% として使ってください．
% 研究の内容によって自分が適切だと思う章構成を考えて下さい．
% もちろん構成について藤原と相談するのも良いと思います．
% \section{テンプレート}
% %サンプルです．必要なければファイルごと削除してください．
% 「はじめに」では研究の背景を述べつつ研究の目的を述べるのが一般的です．
% 仰々しい論文の場合は論文の構成（以降の章構成）を説明したりもします．

% 最初は世の中（といっても情報系分野ぐらいで十分ですが）全体の背景を述べつつ，
% 徐々に的を絞ってどうして自身の研究を行う必要があるのかを上手く説明できると素晴らしいです．
% 研究の背景となる課題や既存研究を紹介する際には必ず参考文献を引用するようにしましょう．

% 「はじめに」は他の章に比べて非常に作成が難しいです．
% どういう流れで書けばいいのか困った際は指導教員に聞くのもよいと思いまし，とりあえず書ける章（自分が実際に手を動かした章）から書いていくのも良いと思います．


% \subsection{論文を書く際の注意事項}
% \begin{itemize}
%     \item 句読点は全角のカンマとピリオドである「，」と「．」を使って下さい．卒論執筆期間中は一時的にIMEの設定を変更して常にカンマとピリオドが使用されるようにしておくと安全だと思います．
%     \item 表記揺れに注意しましょう．
%     \item ですます調はNGです．
%     \item 未定義語は使わない．一般的でない用語が出てくるタイミングで必ず説明を入れる．
% \end{itemize}


\section{序論}
\subsection{研究背景}
ここに研究背景を記述する．研究領域の概要，社会的または学術的な重要性，既存の課題を順に説明する．必要に応じて参考文献を引用する\cite{example2026}．

\subsection{研究目的}
ここに本研究の目的を記述する．研究背景で述べた課題に対して，本研究で何を明らかにするのか，または何を実現するのかを簡潔に述べる．

\subsection{本論文の構成}
本論文の構成を記述する．第2章では研究に必要な前提知識を述べる．第3章では関連研究を整理する．第4章では提案手法を説明する．第5章では実験方法と結果を示す．第6章では結果について考察する．第7章では本論文をまとめる．


\section{準備}
\subsection{用語の定義}
本研究で用いる主要な用語を定義する．

\begin{defi}
ここに定義を記述する．
\end{defi}

\subsection{基礎理論}
研究内容を理解するために必要な基礎理論を記述する．数式は以下のように記述できる．

\begin{equation}
    y = f(x)
    \label{eq:sample}
\end{equation}

\subsection{図表の例}
図を挿入する場合は，画像ファイルを \texttt{images/} に配置し，以下のように記述する．

% \begin{figure}[H]
%     \centering
%     \includegraphics[width=0.8\textwidth]{images/sample.png}
%     \caption{図のキャプション}
%     \label{fig:sample}
% \end{figure}

表は以下のように記述する．

\begin{table}[H]
    \centering
    \caption{表のキャプション}
    \label{tab:sample}
    \begin{tabular}{|c|c|} \hline
        項目 & 値 \\ \hline
        A & 1 \\ \hline
        B & 2 \\ \hline
    \end{tabular}
\end{table}



\section{関連研究}
\subsection{関連研究の概要}
ここに関連研究の概要を記述する．既存研究の目的，手法，得られている知見を整理する．

\subsection{既存手法}
ここに既存手法を記述する．本研究と比較する観点が明確になるように，前提条件，長所，限界を述べる．

\subsection{本研究の位置づけ}
ここに本研究の位置づけを記述する．既存研究との差分と，本研究で扱う課題を明確にする．


\section{提案手法}
\subsection{提案手法の概要}
ここに提案手法の概要を記述する．研究目的に対して，どのような方針で課題を解決するのかを述べる．

\subsection{手法の詳細}
ここに提案手法の詳細を記述する．必要に応じてアルゴリズム，数式，図表を用いて説明する．

\begin{algorithm_step}
    \item[Step 1)] ここに処理の説明を記述する．
    \item[Step 2)] ここに処理の説明を記述する．
    \item[Step 3)] ここに処理の説明を記述する．
\end{algorithm_step}

\subsection{実装}
ここに実装の概要を記述する．コードを掲載する場合は以下のように記述する．

% \begin{lstlisting}[style=python, caption=コード例]
% def main():
%     print("Hello, thesis")
% \end{lstlisting}


\section{実験}
\subsection{実験目的}
ここに実験の目的を記述する．提案手法の何を評価するのかを明確にする．

\subsection{実験条件}
ここに実験条件を記述する．データセット，評価指標，実験環境，比較対象などを整理する．

\subsection{実験結果}
ここに実験結果を記述する．表や図を用いて，結果を客観的に示す．

% ----- 図挿入の使用例（テンプレート参照） -----
% テンプレートファイル: templates/figure_template.tex
% 使い方の例: 以下のコメントを外すか、テンプレートを `\input{templates/figure_template.tex}` で読み込んでください。
%
% 単独図の直接例:
\begin{figure}[htbp]
  \centering
  \includegraphics[width=0.7\linewidth]{images/dummy.png}
  \caption{実験の代表的な結果}
  \label{fig:exp-result}
\end{figure}

% サブ図の例は templates/figure_template.tex を参照してください。


\subsection{結果の分析}
ここに実験結果から読み取れる傾向を記述する．考察章では解釈や原因の議論を中心に扱うため，本節では結果の整理に留める．


\section{考察}
\subsection{結果に対する考察}
ここに実験結果に対する考察を記述する．提案手法が有効であった点，有効でなかった点，その理由を述べる．

\subsection{妥当性}
ここに実験や評価の妥当性を記述する．評価条件の制約や，結果の一般化可能性について述べる．

\subsection{限界}
ここに本研究の限界を記述する．対象範囲，実験条件，手法上の制約を明確にする．


\section{結論}
\subsection{まとめ}
ここに本研究のまとめを記述する．研究目的，提案内容，実験結果，得られた知見を簡潔に整理する．

\subsection{今後の課題}
ここに今後の課題を記述する．本研究で残された課題や，今後取り組むべき発展的な内容を述べる．


% 謝辞
\acknowledgements
ここに謝辞を記述する．指導教員，共同研究者，研究室のメンバー，家族など，研究の遂行にあたって支援を受けた相手に感謝を述べる．


% 参考文献
% BibTeX を使って参考文献のリストをつくって下さい．
% proceedings.bib: 雑誌名や国際会議の名称をマクロ化しています
% ref.bib: ここに参考文献を書くと良いと思います．元々博士論文に使った論文の情報が載っています．
% ref.bib以外のファイルに分けることも可能です．単純にカンマで区切ってファイル名（.bib除く）を
% \bibliographyの引数に追加してください．
\bibliography{proceedings,ref}

% 付録
% 付録には，本文に直接載せるべきではない，文字通り付録となる要素を載せます．
% 例えば，詳細な数式の展開であったり，アンケートや実験の詳細な（ほぼ完全な）結果だったりを載せます．
% 必要であれば以下の3行のコメントアウトを外して付録を加えてください．
\appendix
\section{付録のタイトル}
ここに付録を記述する．本文に直接載せると冗長になる詳細な数式，追加の実験結果，アンケート項目などを配置する．

\section{実験コード}
本実験で使用したコードを以下に示す．

\subsection{メインスクリプト}
以下は実験で使用したPythonスクリプト（\texttt{src/main.py}）である．このコードは基本的な実装例として提供される．

\lstinputlisting[language=Python, caption={main.py - メインスクリプト}, label={code:main}]{src/main.py}

% 複数のコードを挿入する場合の例:
% \subsection{別のスクリプト}
% \lstinputlisting[language=Java, caption={別の言語の例}]{src/example.java}
%
% 参考: このテンプレートの詳細は templates/code\_template.tex を参照してください．


\end{document}
